% ============================================================
% Chapter 6: Network Security: Attacks, Defences and Botnet Topologies
% Language: Greek — edit freely; regenerate from src/content if preferred.
% ============================================================
\chapter{Ασφάλεια Δικτύων: Επιθέσεις, Άμυνες και Τοπολογίες Botnet}\label{ch:6}
\chaptermeta{DoS και DDoS · Ενίσχυση · Botnet και διοίκηση-έλεγχος · Επιθέσεις ενδιάμεσου · Τείχη προστασίας · VPN · Κατάτμηση, NAC και ασφάλεια ασύρματων δικτύων}{Επίπεδο: Μεσαίο επίπεδο · Δίκτυα}{Χρόνος μελέτης: 8–10 ώρες}

Τα δίκτυα συνδέουν τα πάντα, και ό,τι συνδέεται μπορεί επίσης να δεχθεί επίθεση. Το κεφάλαιο εξετάζει πώς οι αντίπαλοι κάνουν κατάχρηση των δικτυακών πρωτοκόλλων για να διαταράξουν τη διαθεσιμότητα, να υποκλέψουν ή να αλλοιώσουν την κίνηση και να συντονίσουν χιλιάδες παραβιασμένες συσκευές. Προϋποθέτει βασική εξοικείωση με το μοντέλο TCP/IP· όπου μια λεπτομέρεια πρωτοκόλλου είναι απαραίτητη, εξηγείται στο κείμενο.

Το πρώτο μισό του κεφαλαίου εστιάζει στις επιθέσεις: την άρνηση υπηρεσίας και τις κατανεμημένες και ενισχυμένες μορφές της, την αρχιτεκτονική και τον κύκλο ζωής των botnet, καθώς και τις τεχνικές ενδιάμεσου (man-in-the-middle). Το δεύτερο μισό στρέφεται στην άμυνα. Δείχνει πώς λειτουργεί στην πράξη η πολυεπίπεδη αντιμετώπιση επιθέσεων DDoS και παρουσιάζει τα δομικά στοιχεία της δικτυακής άμυνας —τείχη προστασίας, εικονικά ιδιωτικά δίκτυα, κατάτμηση, έλεγχο πρόσβασης στο δίκτυο και ασφάλεια ασύρματων δικτύων— με σαφή διατύπωση του τι μπορεί και τι δεν μπορεί να επιτύχει το καθένα.

\begin{outcomesbox}{Μαθησιακά αποτελέσματα}
Μετά τη μελέτη του κεφαλαίου θα πρέπει να μπορείτε:
\begin{itemize}
  \item Να διακρίνετε τις επιθέσεις άρνησης υπηρεσίας ογκομετρικού τύπου, επιπέδου πρωτοκόλλου και επιπέδου εφαρμογής και να εξηγείτε την ενίσχυση (amplification).
  \item Να συγκρίνετε τις κεντρικές και τις ομότιμες αρχιτεκτονικές botnet και να περιγράφετε τον κύκλο ζωής ενός botnet.
  \item Να εξηγείτε την πλαστογράφηση ARP, την πλαστογράφηση DNS και την απογύμνωση TLS, καθώς και τα μέτρα που τις αποτρέπουν.
  \item Να σχεδιάζετε μια πολυεπίπεδη άμυνα κατά DDoS, από την άκρη του δικτύου έως την εφαρμογή.
  \item Να περιγράφετε τον ρόλο και τα όρια των τειχών προστασίας, των VPN, της κατάτμησης, του 802.1X και του WPA3.
\end{itemize}
\end{outcomesbox}

\section{Άρνηση υπηρεσίας: DoS, DDoS και ενίσχυση}\label{sec:6.1}
Μια επίθεση \textbf{άρνησης υπηρεσίας} (Denial of Service, DoS) στοχεύει να καταστήσει μια υπηρεσία μη διαθέσιμη στους νόμιμους χρήστες της. Όταν η κίνηση της επίθεσης προέρχεται ταυτόχρονα από πολλές πηγές, ονομάζεται \textbf{κατανεμημένη άρνηση υπηρεσίας} (Distributed Denial of Service, DDoS) και γίνεται πολύ δυσκολότερο να αποκλειστεί, επειδή δεν υπάρχει μία πηγή που να μπορεί απλώς να φιλτραριστεί. Οι επιθέσεις DDoS χωρίζονται σε τρεις γενικές οικογένειες. Οι \textbf{ογκομετρικές} κατακλύζουν το εύρος ζώνης του θύματος. Οι επιθέσεις \textbf{επιπέδου πρωτοκόλλου} εξαντλούν τους πόρους κατάστασης διακομιστών ή δικτυακών συσκευών· η κλασική \emph{πλημμύρα SYN} στέλνει πολλά αιτήματα σύνδεσης TCP χωρίς να ολοκληρώνει τη χειραψία, γεμίζοντας τον πίνακα ημιανοιχτών συνδέσεων του διακομιστή. Οι επιθέσεις \textbf{επιπέδου εφαρμογής} στέλνουν φαινομενικά νόμιμα αλλά απαιτητικά αιτήματα —όπως ερωτήματα αναζήτησης— που εξαντλούν την ίδια την εφαρμογή.

Η \textbf{ενίσχυση} (amplification) επιτρέπει σε έναν επιτιθέμενο με περιορισμένο εύρος ζώνης να παράγει τεράστιες πλημμύρες κίνησης. Ο επιτιθέμενος στέλνει μικρά αιτήματα με πλαστογραφημένη διεύθυνση προέλευσης (τη διεύθυνση του θύματος) σε δημόσιους διακομιστές, οι οποίοι επιστρέφουν πολύ μεγαλύτερες απαντήσεις. Πρωτόκολλα όπως τα DNS, NTP και memcached έχουν χρησιμοποιηθεί με αυτόν τον τρόπο, με συντελεστές ενίσχυσης από μερικές δεκάδες έως δεκάδες χιλιάδες. Δύο μέτρα αντιμετωπίζουν τη ρίζα του προβλήματος: τα δίκτυα πρέπει να απορρίπτουν εξερχόμενα πακέτα με πλαστογραφημένες διευθύνσεις προέλευσης (\emph{φιλτράρισμα εισόδου}, BCP 38) και οι διαχειριστές δεν πρέπει να εκθέτουν στο διαδίκτυο ανοιχτούς αναλυτές DNS ή άλλους ενισχυτές.

\section{Botnet: αρχιτεκτονικές, κύκλος ζωής και στρατολόγηση συσκευών IoT}\label{sec:6.2}
Ένα \textbf{botnet} είναι ένα δίκτυο παραβιασμένων συσκευών («bots») που ελέγχονται απομακρυσμένα από έναν διαχειριστή. Σε μια \textbf{κεντρική} αρχιτεκτονική, τα bots λαμβάνουν εντολές από έναν ή λίγους διακομιστές διοίκησης και ελέγχου (Command and Control, C2), ιστορικά μέσω IRC και σήμερα κυρίως μέσω HTTPS. Ο σχεδιασμός αυτός είναι απλός αλλά εύθραυστος: η κατάσχεση των διακομιστών C2 απενεργοποιεί το botnet. Τα \textbf{ομότιμα} (peer-to-peer) botnet κατανέμουν τις εντολές στα ίδια τα bots, κάτι που καθιστά την εξάρθρωσή τους πολύ δυσκολότερη. Οι διαχειριστές χρησιμοποιούν επίσης \emph{αλγορίθμους παραγωγής ονομάτων τομέα} (DGA), οι οποίοι δημιουργούν χιλιάδες υποψήφια ονόματα κάθε ημέρα, ώστε οι αμυνόμενοι να μην μπορούν να τα αποκλείσουν όλα εκ των προτέρων.

Ο κύκλος ζωής ενός botnet περιλαμβάνει τη \textbf{στρατολόγηση} (μόλυνση νέων συσκευών), τη \textbf{διοίκηση και τον έλεγχο}, την \textbf{αξιοποίηση} (DDoS επί πληρωμή, ανεπιθύμητη αλληλογραφία, μαζική δοκιμή διαπιστευτηρίων, εξόρυξη κρυπτονομισμάτων) και τη \textbf{συντήρηση} (ενημέρωση και προστασία των bots). Το botnet \textbf{Mirai} του 2016 έδειξε πόσο αδύναμο ήταν το Διαδίκτυο των Πραγμάτων: απλώς δοκίμαζε μια λίστα περίπου εξήντα εργοστασιακών ονομάτων χρήστη και κωδικών σε κάμερες και δρομολογητές, στρατολόγησε εκατοντάδες χιλιάδες συσκευές και εξαπέλυσε επιθέσεις που διατάραξαν σημαντικές υπηρεσίες του διαδικτύου. Το δίδαγμα για κατασκευαστές και διαχειριστές είναι σαφές —μοναδικά διαπιστευτήρια, αυτόματες ενημερώσεις και καμία περιττή απομακρυσμένη υπηρεσία.

\section{Υποκλοπή και αλλοίωση: επιθέσεις ενδιάμεσου}\label{sec:6.3}
Σε μια \textbf{επίθεση ενδιάμεσου} (Man-in-the-Middle, MITM), ο αντίπαλος παρεμβάλλεται ανάμεσα σε δύο μέρη που επικοινωνούν, ώστε η κίνηση να διέρχεται από ένα σύστημα που ελέγχει. Σε ένα τοπικό δίκτυο, η \textbf{πλαστογράφηση ARP} το επιτυγχάνει αποστέλλοντας ψευδή μηνύματα που συσχετίζουν τη φυσική διεύθυνση του επιτιθέμενου με τη διεύθυνση IP της προεπιλεγμένης πύλης· τα θύματα στέλνουν στη συνέχεια την κίνησή τους στον επιτιθέμενο. Η \textbf{πλαστογράφηση DNS} επιστρέφει ψευδείς απαντήσεις σε ερωτήματα ονομάτων, ανακατευθύνοντας τους χρήστες σε κακόβουλους διακομιστές. Σε δημόσια δίκτυα Wi-Fi, ένα σημείο πρόσβασης \emph{«κακός δίδυμος»} (evil twin) με γνώριμο όνομα δικτύου μπορεί να προσελκύσει ανυποψίαστους χρήστες.

Μόλις βρεθεί στη μέση, ο επιτιθέμενος μπορεί είτε να \textbf{υποκλέπτει} παθητικά τη μη κρυπτογραφημένη κίνηση είτε να την τροποποιεί ενεργά —για παράδειγμα μέσω της \emph{απογύμνωσης TLS} (TLS stripping), η οποία υποβαθμίζει μια σύνδεση από HTTPS σε HTTP. Η κρυπτογραφία είναι το θεμελιώδες αντίμετρο: το σωστά επικυρωμένο TLS (Κεφάλαιο 7) καθιστά την υποκλαπείσα κίνηση μη αναγνώσιμη και την αλλοίωση ανιχνεύσιμη, ενώ ο μηχανισμός \textbf{HTTP Strict Transport Security (HSTS)} αποτρέπει την υποβάθμιση. Σε επίπεδο δικτύου, η \emph{δυναμική επιθεώρηση ARP} και η \emph{παρακολούθηση DHCP} (DHCP snooping) στους μεταγωγείς, η επικύρωση DNSSEC και ο έλεγχος ταυτότητας 802.1X περιορίζουν τις ευκαιρίες υποκλοπής.

\section{Πολυεπίπεδη άμυνα κατά DDoS στην πράξη}\label{sec:6.4}
Καμία μεμονωμένη συσκευή δεν μπορεί να απορροφήσει κάθε επίθεση DDoS· γι' αυτό η άμυνα οργανώνεται σε επίπεδα, καθένα από τα οποία χειρίζεται αυτό που κάνει καλύτερα. Στο \textbf{ανάντη} επίπεδο, οι πάροχοι διαδικτύου και οι υπηρεσίες «καθαρισμού» κίνησης στο νέφος απορροφούν τις ογκομετρικές πλημμύρες, αξιοποιώντας χωρητικότητα πολύ μεγαλύτερη από τη σύνδεση οποιουδήποτε μεμονωμένου οργανισμού· η δρομολόγηση anycast κατανέμει το φορτίο σε πολλά κέντρα δεδομένων. Στην \textbf{άκρη του δικτύου}, οι δρομολογητές εφαρμόζουν λίστες πρόσβασης και όρια ρυθμού, ενώ η \emph{απομακρυσμένα ενεργοποιούμενη «μαύρη τρύπα»} (remote-triggered black-holing) μπορεί να θυσιάσει μία διεύθυνση που δέχεται επίθεση για να προστατεύσει το υπόλοιπο δίκτυο.

Πιο κοντά στην υπηρεσία, \textbf{άμυνες επιπέδου πρωτοκόλλου}, όπως τα \emph{SYN cookies}, επιτρέπουν στους διακομιστές να χειρίζονται πλημμύρες ημιανοιχτών συνδέσεων χωρίς να αποθηκεύουν κατάσταση. Στο επίπεδο της \textbf{εφαρμογής}, τα δίκτυα διανομής περιεχομένου (CDN), τα τείχη προστασίας εφαρμογών ιστού (WAF), η προσωρινή αποθήκευση και ο περιορισμός ρυθμού ανά πελάτη φιλτράρουν τα απαιτητικά αιτήματα, ενώ δοκιμασίες όπως τα CAPTCHA διαχωρίζουν τους ανθρώπους από τα bots. Τέλος, ένα τεκμηριωμένο \textbf{εγχειρίδιο ενεργειών} (runbook) —ποιον καλούμε, πώς ενεργοποιούμε τον καθαρισμό κίνησης, πώς επικοινωνούμε με τους πελάτες— εξασφαλίζει ότι ο οργανισμός αντιδρά σε λεπτά και όχι σε ώρες.

\section{Τείχη προστασίας και εικονικά ιδιωτικά δίκτυα}\label{sec:6.5}
Ένα \textbf{τείχος προστασίας} (firewall) επιβάλλει μια πολιτική σχετικά με το ποια κίνηση επιτρέπεται να διέρχεται μεταξύ ζωνών του δικτύου. Τα \emph{φίλτρα πακέτων} εξετάζουν μεμονωμένα πακέτα με βάση τη διεύθυνση, τη θύρα και το πρωτόκολλο. Τα τείχη προστασίας \emph{με διατήρηση κατάστασης} (stateful) παρακολουθούν τις συνδέσεις, ώστε μια απάντηση να επιτρέπεται μόνο αν αντιστοιχεί σε αίτημα που έχει σταλεί. Τα \emph{τείχη προστασίας νέας γενιάς} αναγνωρίζουν επιπλέον εφαρμογές και χρήστες και μπορεί να ενσωματώνουν πρόληψη εισβολών. Ανεξάρτητα από τη γενιά, ένας κανόνας σχεδίασης είναι θεμελιώδης: η πολιτική πρέπει να \textbf{απορρίπτει εξ ορισμού} και να επιτρέπει μόνο τις ρητά αναγκαίες ροές, σύμφωνα με την αρχή των ασφαλών προεπιλογών.

Ένα \textbf{εικονικό ιδιωτικό δίκτυο} (Virtual Private Network, VPN) δημιουργεί ένα κρυπτογραφημένο τούνελ μέσα από ένα μη αξιόπιστο δίκτυο. Το \textbf{IPsec} λειτουργεί στο επίπεδο δικτύου και χρησιμοποιείται συχνά για συνδέσεις μεταξύ εγκαταστάσεων· τα \textbf{VPN που βασίζονται σε TLS} είναι βολικά για απομακρυσμένους χρήστες· το \textbf{WireGuard} είναι ένα σύγχρονο πρωτόκολλο με σκόπιμα μικρό όγκο κώδικα και σταθερή, σύγχρονη κρυπτογραφία. Είναι σημαντικό να κατανοηθεί τι \emph{δεν} κάνει ένα VPN: προστατεύει την κίνηση κατά τη μεταφορά, αλλά δεν καθιστά αξιόπιστη τη συσκευή που συνδέεται. Ένας παραβιασμένος φορητός υπολογιστής που συνδέεται μέσω VPN φέρνει τον επιτιθέμενο μέσα στο δίκτυο —ένας από τους λόγους για τους οποίους οι οργανισμοί μεταβαίνουν σε πρόσβαση μηδενικής εμπιστοσύνης (Κεφάλαιο 13).

\section{Κατάτμηση, έλεγχος πρόσβασης στο δίκτυο και ασφάλεια ασύρματων δικτύων}\label{sec:6.6}
Η \textbf{κατάτμηση δικτύου} (segmentation) διαιρεί ένα δίκτυο σε ζώνες —για παράδειγμα σταθμούς εργασίας χρηστών, διακομιστές, διεπαφές διαχείρισης και Wi-Fi επισκεπτών— και επιτρέπει μόνο την κίνηση που πραγματικά χρειάζεται κάθε ζώνη. Η κατάτμηση δεν αποτρέπει την πρώτη παραβίαση, περιορίζει όμως το πόσο μακριά μπορεί να κινηθεί ο επιτιθέμενος στη συνέχεια· η περίπτωση του WannaCry στο Κεφάλαιο 4 έδειξε το κόστος των επίπεδων δικτύων. Η \emph{μικροκατάτμηση} εφαρμόζει την ίδια ιδέα έως το επίπεδο των μεμονωμένων φόρτων εργασίας.

Ο \textbf{έλεγχος πρόσβασης στο δίκτυο} (Network Access Control, NAC) αποφασίζει αν μια συσκευή επιτρέπεται καν να συνδεθεί στο δίκτυο. Το πρότυπο \textbf{IEEE 802.1X} απαιτεί από κάθε συσκευή να επαληθεύσει την ταυτότητά της —συνήθως με πιστοποιητικό— πριν ενεργοποιηθεί η θύρα του μεταγωγέα ή η ασύρματη σύνδεση, και μπορεί να τοποθετεί μη συμμορφούμενες συσκευές σε δίκτυο καραντίνας. Η ασφάλεια των ασύρματων δικτύων έχει εξελιχθεί σημαντικά: το WEP είναι πλήρως παραβιασμένο, το WPA2 είναι ευάλωτο σε μαντεψιά κωδικών εκτός σύνδεσης όταν χρησιμοποιούνται αδύναμες φράσεις πρόσβασης, ενώ το \textbf{WPA3} εισάγει τη χειραψία SAE, η οποία αντιστέκεται σε τέτοιες μαντεψιές και παρέχει εμπρόσθια μυστικότητα. Στις επιχειρήσεις, η συνιστώμενη ρύθμιση παραμένει το WPA2/WPA3-Enterprise με 802.1X.

\section*{Βασικοί όροι}
\begin{description}[style=nextline,leftmargin=1.2em]
  \item[DDoS] Κατανεμημένη άρνηση υπηρεσίας: επίθεση κατά της διαθεσιμότητας από πολλές πηγές ταυτόχρονα.
  \item[Ενίσχυση] Αξιοποίηση τρίτων διακομιστών που επιστρέφουν μεγάλες απαντήσεις σε μικρά, πλαστογραφημένα αιτήματα.
  \item[Botnet / C2] Δίκτυο παραβιασμένων συσκευών και το κανάλι διοίκησης και ελέγχου που τις κατευθύνει.
  \item[Πλαστογράφηση ARP] Αποστολή ψευδών μηνυμάτων ARP για την ανακατεύθυνση της τοπικής κίνησης μέσω του επιτιθέμενου.
  \item[Τείχος προστασίας με κατάσταση] Τείχος προστασίας που παρακολουθεί τις συνδέσεις και επιτρέπει απαντήσεις μόνο σε νόμιμα αιτήματα.
  \item[802.1X] Έλεγχος πρόσβασης ανά θύρα που επαληθεύει την ταυτότητα των συσκευών πριν τους παραχωρήσει συνδεσιμότητα.
\end{description}

\section*{Σύνοψη κεφαλαίου}
\begin{itemize}
  \item Οι επιθέσεις DDoS στοχεύουν το εύρος ζώνης, την κατάσταση των πρωτοκόλλων ή τους πόρους της εφαρμογής· η ενίσχυση πολλαπλασιάζει τη δύναμη του επιτιθέμενου μέσω πλαστογράφησης.
  \item Τα botnet ελέγχονται μέσω κεντρικής ή ομότιμης διοίκησης-ελέγχου και αξιοποιούνται με πολλούς τρόπους· το Mirai ανέδειξε την αδυναμία των εργοστασιακών διαπιστευτηρίων στις συσκευές IoT.
  \item Οι επιθέσεις ενδιάμεσου εκμεταλλεύονται την εμπιστοσύνη στα τοπικά πρωτόκολλα· τα κύρια αντίμετρα είναι το επικυρωμένο TLS, το HSTS και οι προστασίες στους μεταγωγείς.
  \item Η άμυνα κατά DDoS είναι πολυεπίπεδη, από τον καθαρισμό κίνησης στο ανάντη δίκτυο έως τον περιορισμό ρυθμού στην εφαρμογή, και εξαρτάται από ένα προετοιμασμένο εγχειρίδιο ενεργειών.
  \item Τα τείχη προστασίας, τα VPN, η κατάτμηση και ο NAC επιλύουν συγκεκριμένα προβλήματα —κανένα δεν καθιστά αξιόπιστη μια παραβιασμένη συσκευή.
\end{itemize}

\section*{Ερωτήσεις ανασκόπησης}
\begin{enumerate}
  \item Εξηγήστε γιατί ένας επιτιθέμενος που χρησιμοποιεί ενίσχυση DNS πρέπει να πλαστογραφεί τη διεύθυνση IP προέλευσης των αιτημάτων του.
  \item Γιατί τα ομότιμα botnet εξαρθρώνονται δυσκολότερα από τα κεντρικά;
  \item Περιγράψτε πώς το HSTS αποτρέπει μια επίθεση απογύμνωσης TLS σε δημόσιο δίκτυο Wi-Fi.
  \item Μια εταιρεία δηλώνει: «Όλοι οι απομακρυσμένοι εργαζόμενοι χρησιμοποιούν το VPN μας, άρα οι φορητοί τους υπολογιστές είναι ασφαλείς». Αξιολογήστε κριτικά αυτή τη δήλωση.
\end{enumerate}
